\introduction

La infraestructura de telecomunicaciones es un pilar fundamental para el desarrollo y la prestación
de servicios digitales en la era contemporánea. En el contexto nacional \citep{Press2011, Hernandez2021},
la red de Internet presenta insuficiencias en su infraestructura, lo que se traduce de manera tangible
en bajas velocidades de conexión para los usuarios finales. Esta limitación del entorno tecnológico
impacta directamente en la calidad y la eficiencia de los servicios en línea que operan en el país,
generando desafíos técnicos que requieren soluciones innovadoras y adaptadas a estas condiciones
particulares.

Entre estos servicios se encuentra Picta \citep{UCI2025}, la principal plataforma de
streaming nacional, la cual se ve significativamente afectada por esta problemática. Su operación
depende de la capacidad de más de 700 canales activos para subir y publicar su contenido de video de
forma fluida. Sin embargo, los procesos de subida y renderizado de videos, inherentemente intensivos
en el consumo de ancho de banda y recursos computacionales del servidor, se ven severamente
obstaculizados por las limitaciones de la red. Esto se manifiesta en tiempos de carga excesivamente
largos, interrupciones y una alta demanda sostenida sobre los servidores de la plataforma. La
consecuencia directa es una experiencia de usuario deficiente para los creadores de contenido y una
barrera para la escalabilidad y confiabilidad del servicio de Picta, poniendo en riesgo su crecimiento
sostenible.

Partiendo de esta problemática, se formula el \textbf{problema científico}: ¿Cómo contribuir al proceso 
de subida y renderizado de videos en la plataforma Picta para mitigar los efectos de las limitaciones de 
la infraestructura de red nacional, garantizando una experiencia de usuario fluida y uso eficiente de los 
recursos del servidor? Para abordar este problema, se define como \textbf{objeto de estudio} el proceso 
de subida y renderizado de videos, y como \textbf{campo de acción} específico, este mismo proceso dentro 
del ecosistema de la plataforma Picta.

El \textbf{objetivo general} de esta investigación es desarrollar un sistema de subida fragmentada y 
renderizado asíncrono de videos para la plataforma Picta que mejore la eficiencia en el uso de recursos 
y la experiencia del usuario frente a las limitaciones de la infraestructura de red nacional. Para alcanzar 
este propósito, se establecen los siguientes \textbf{objetivos específicos}:

\begin{enumerate}
  \item Caracterizar el flujo actual de subida y renderizado de videos en Picta para identificar
  cuellos de botella y oportunidades de mejora.

  \item Analizar aplicaciones con funcionalidades similares a las requeridas.

  \item Diseñar un sistema de subida fragmentada que permita la persistencia del progreso y el
  control del usuario durante la carga de videos.

  \item Implementar un mecanismo de renderizado asíncrono que separe la subida del video de su
  procesamiento final.

  \item Evaluar el rendimiento del sistema implementado en términos de tiempo de subida, uso de
  ancho de banda y experiencia del usuario.
\end{enumerate}

Para alcanzar los objetivos propuestos, se emplean \textbf{métodos teóricos y empíricos} que 
guían el proceso investigativo. Entre los \textbf{métodos teóricos} se aplican:

\begin{itemize}
  \item \textbf{Análisis-síntesis:} Permite descomponer el proceso de subida y renderizado en 
  sus componentes fundamentales ---transferencia, validación, transcoding--- para identificar 
  cuellos de botella específicos, y posteriormente sintetizar una solución integrada mediante la 
  arquitectura dirigida por eventos.

  \item \textbf{Histórico-lógico:} Posibilita estudiar la trayectoria histórica de los 
  protocolos de transferencia de archivos y sistemas de almacenamiento, extrayendo los principios 
  lógicos fundamentales que se aplican en el diseño de la solución propuesta.

  \item \textbf{Modelación:} Facilita la representación abstracta del sistema mediante diagramas 
  de flujo, patrones arquitectónicos y de diseño, permitiendo estructurar la solución antes de su 
  implementación física.
\end{itemize}

En cuanto a los \textbf{métodos empíricos}, se emplean:

\begin{itemize}
  \item \textbf{Estudio comparativo:} Se realiza un análisis de sistemas homólogos con 
  funcionalidades similares para identificar fortalezas y limitaciones que informan el diseño de 
  la solución especializada para Picta.

  \item \textbf{Observación:} Se aplica para caracterizar el flujo actual de subida y 
  renderizado en la plataforma, identificando patrones de fallo y limitaciones derivadas de la 
  infraestructura de red nacional.

  \item \textbf{Experimentación:} Se ejecutan pruebas de rendimiento del sistema implementado 
  para validar mejoras en tiempo de subida, uso de ancho de banda y experiencia del usuario, 
  comparando métricas antes y después de la intervención.
\end{itemize}

Este documento está organizado en tres capítulos, conclusiones generales y bibliografía. Los
capítulos se estructuran de la siguiente forma:

\begin{enumerate}
  \item \textbf{Capítulo 1: Fundamentación teórica.} Este capítulo realiza un estudio de soluciones
  homólogas, del flujo actual de subida y renderizado de videos en Picta para identificar cuellos de
  botella y oportunidades de mejora. Para este propósito, además se definen herramientas, técnicas,
  tecnologías y metodologías que se emplearán para alcanzar el objetivo general.

  \item \textbf{Capítulo 2: Diseño de la aplicación.} Se realiza el proceso de ingeniería
  correspondiente a la metodología escogida. Se definen los requisitos que debería cumplir el
  programa en la forma de historias de usuario, se presentan los artefactos generados por la
  metodología de desarrollo y se concluye mencionando los patrones arquitectónicos y de diseño que
  sigue la aplicación.

  \item \textbf{Capítulo 3: Implementación y pruebas.} Se definen tareas de ingeniería y se realiza
  el proceso de pruebas de la aplicación. Se finaliza realizando una validación de la solución
  propuesta.
\end{enumerate}

La relevancia científica de esta propuesta radica en la transición hacia un modelo de procesamiento
asíncrono y distribuido, rompiendo con el esquema tradicional de transferencias monolíticas que
penalizan al usuario ante cualquier micro-corte de red. Al emplear una arquitectura dirigida por
eventos, se busca transformar la infraestructura lógica de la plataforma para que sea capaz de
gestionar cargas de trabajo de forma elástica y resiliente \citep{Richards2020}. Este enfoque
permite que el sistema reaccione de manera dinámica ante la disponibilidad de recursos, asegurando
que la etapa de renderización ---crítica por su consumo de CPU--- no comprometa la disponibilidad
inmediata de la interfaz de usuario ni la integridad de los datos en tránsito.

Desde una perspectiva práctica, este trabajo se inserta en el esfuerzo por fortalecer la soberanía
tecnológica y mejorar la accesibilidad de los servicios digitales en el contexto nacional
\citep{Fuentes2025}. Al optimizar el uso de la red mediante la fragmentación de paquetes, se
democratiza la creación de contenido, permitiendo que usuarios con conexiones inestables o de bajo
ancho de banda puedan contribuir de manera efectiva al ecosistema audiovisual de Picta. De este
modo, la solución propuesta trasciende la optimización técnica de un módulo específico para
convertirse en un habilitador de escalabilidad que garantiza la continuidad del servicio frente a
una demanda creciente de tráfico y almacenamiento \citep{Avizienis2004}.

Como \textbf{resultados esperados}, se prevé una mejora tangible en la experiencia de los usuarios
que suben contenido, otorgándoles un mayor control sobre el proceso; una reducción significativa en
el consumo de ancho de banda durante el proceso; y una optimización en la utilización de los recursos
del servidor, lo que en conjunto contribuirá a una mayor robustez y escalabilidad de la plataforma
Picta \citep{Pressman2021}. La presente tesis se organiza para documentar de manera exhaustiva este
proceso, desde el análisis inicial hasta la implementación y validación de la solución propuesta,
con el fin de ofrecer una respuesta técnica viable a un problema de relevancia en el contexto digital
nacional.